\section{Introduction}

Federated learning lets multiple clients train collaboratively without sharing their private training data. In model-heterogeneous federated learning (HFL), clients may also use different architectures, which rules out direct parameter averaging and motivates knowledge transfer through predictions, prototypes, or public data. FedDF distills ensemble predictions over unlabeled data \citep{lin2020feddf}; RHFL addresses noisy labels under model heterogeneity \citep{fang2022rhfl}; and AugHFL and its TPAMI extension RAHFL study heterogeneous collaboration under common data corruption \citep{fang2023aughfl,fang2025rahfl}. These methods treat corruption mainly as label-preserving nuisance variation, consistent with standard corruption benchmarks \citep{hendrycks2019corruptions,cubuk2020augmix}.

We study a different failure mode. Suppose a client sees most images of one class with blur and most images of another class with noise, while another client has a different class--corruption association. Corruption can then become predictive of the task label. We call this controlled setting \textbf{class--corruption entanglement in HFL (CLE-HFL)}. The central question is not only whether corruption reduces accuracy, but whether changing the corruption operator for the same semantic source moves predictions toward the classes associated with that operator during training. Ordinary corrupted accuracy cannot separate this directional shortcut from image degradation, label imbalance, model capacity, or non-IID optimization.

We therefore use a source-paired operator diagnostic. AugMix-style Jensen--Shannon consistency can align augmentations around one corrupted view while leaving cross-operator class routing unchanged \citep{cubuk2020augmix}. A pseudo-environment score can also improve even when the model still follows the true class--corruption binding. We define \textbf{Directional Shortcut Alignment (DSA)} as the change in probability mass toward the classes pre-registered as bound to the applied operator, using the same semantic source across operator interventions. DSA measures response along the training binding under assumptions of semantic preservation, valid operator intervention, pre-registered bindings, source pairing, and evaluation isolation. We also show that, with an unconstrained proxy error channel, proxy statistics alone do not identify true CLE dependence; final mitigation claims therefore require a target-aligned behavioral endpoint.

The paired diagnostic also lets us ask where the shortcut forms. In a seed-0 matched HFL-versus-Local factorial, the strong-CLE contrast is similar with and without communication, while the additional communication contrast is much smaller. We describe this pattern as \textbf{predominantly local-first} and target the client objective rather than redesigning communication. Our intervention uses a \textbf{coarse family-level Public Environment Witness (PEW)} trained on public synthetic corruptions and \textbf{Balanced Environment Risk (BER)}, which compresses support-induced risk-mass differences across pseudo-environments within each task class. PEW is explicitly taxonomy-assisted and does not read private true corruption operators or families.

We test the intervention along two comparison axes. Axis I compares ERM, a protocol-matched CVaR-DRO control, PEW+GroupDRO with the same PEW grouping information, and PEW+BER on a held-out binding map over training seeds 0/1/2. PEW+BER has the lowest pooled DSA and the highest pooled task utility among these four matched objectives. Shared validation changes only the binding map or the HFL communication base; DSA suppression reproduces in both tests, although the FedDF-fidelity utility gate fails. Axis II is reserved for protocol-matched HFL mechanism-family context and remains pending until Formal results are available. Our contributions are:

\begin{enumerate}
\item \textbf{CLE-HFL problem formulation.} We formulate client-specific class--corruption bindings in model-heterogeneous FL and separate directional shortcut use from generic corruption difficulty.
\item \textbf{Paired diagnosis and identification boundary.} We introduce source-paired operator evaluation and \textbf{Directional Shortcut Alignment (DSA)}, characterize its identification and source-level inference, and prove a proxy non-identifiability result that separates proxy improvement from target-aligned CLE mitigation.
\item \textbf{Local-versus-communication formation analysis.} A matched HFL-versus-Local factorial shows a predominantly local-first pattern in the evaluated seed-0 setting, with a smaller pooled communication add-on.
\item \textbf{Taxonomy-assisted local mitigation and controlled validation.} We combine coarse family-level PEW with BER and evaluate it against matched ERM, tail risk, and group-robustness controls, then test different binding maps and communication bases while reporting utility and taxonomy limits separately.
\end{enumerate}

Prior work already studies federated spurious features, invariant or group-robust learning, environment inference, and counterfactual shortcut evaluation \citep{wang2024pflspurious,tang2024fedpin,ma2024fedcd,creager2021eiil,sagawa2020groupdro,vigneshwaran2026counterfactual}. Our scope is their specific intersection with client-dependent class--corruption bindings in model-heterogeneous FL, together with a paired directional diagnostic and local-versus-communication formation analysis.
